\documentclass[../../main.tex]{subfiles}
\graphicspath{{../../image/}{../../figure/}}
\begin{document}

\section{古典解的崩溃}\label{section:ch04-经典解的崩溃}

这里我们将结识两种不同的机制，即“波破碎”与“质量囚禁”，它们可能导致非线性双曲守恒律组 Cauchy 问题古典解的崩溃。

我们首先看到，非线性迫使压缩波剖面变陡并最终破碎，从而使解的一个导数爆破。这将在单个空间变量中最简单的非线性双曲守恒律--Burgers 方程--的框架内得到演示：
\begin{align}
\partial_{t}u(x,t)+\partial_{x}\left(\tfrac{1}{2}u^{2}(x,t)\right)=0.
\label{align:ch04-Burgers}
\end{align}
这个看似简单的方程贯穿于双曲守恒律理论，因为它在一般方程组的分析中反复、自发地出现；例如见\cref{chapter:ch07-一维双曲平衡律组}。


设 $u(x,t)$ 是\eqref{align:ch04-Burgers}以初值 $u_{0}(\cdot)$ 在某时间区间 $[0,T)$ 上的古典解。与该解相联系的\eqref{align:ch04-Burgers}的特征线是常微分方程 $\mathrm{d}x/\mathrm{d}t=u(x,t)$ 的积分曲线。记对时间的微分为点号，$\cdot=\partial_{t}+u\partial_{x}$ 表示沿特征方向的导数，则可把\eqref{align:ch04-Burgers}改写为 $\dot{u}=0$，这表明 $u$ 沿特征线为常数。这反过来又意味着特征线是直线。

令 $\partial_{x}u=v$，并对 $x$ 微分\eqref{align:ch04-Burgers}，得方程 $\partial_{t}v+u\partial_{x}v+v^{2}=0$，即 $\dot{v}+v^{2}=0$。因此，从任一点 $(\bar{x},0)$（其中 $u_{0}'(\bar{x})<0$）发出的特征线上，$|\partial_{x}u|$ 是一个递增函数，并在时刻 $\bar{t}=[-u_{0}'(\bar{x})]^{-1}$ 处爆破。于是显然，$u(x,t)$ 作为古典解必在时刻 $\bar{t}$ 或之前崩溃。

为获得关于波破碎的另一富有启发性的视角，让我们对上述解定义域上的任一点 $(x,t)$ 关联数 $y=x-tu(x,t)$。于是 $y$ 标记了通过点 $(x,t)$ 的、与 $u$ 相联系的特征线在 $x$-轴上的截距。映射 $(x,t)\mapsto(y,t)$ 诱导出一个新的坐标系。在连续物理的精神下，可把 $(x,t)$ 视为“Euler 坐标”，而把 $(y,t)$ 视为“Lagrange 坐标”。在 Lagrange 坐标下，解取简单形式 $u(y,t)=u_{0}(y)$，导数有界：$\partial_{y}u(y,t)=u_{0}'(y)$，$\partial_{t}u(y,t)=0$。然而，当换回 Euler 坐标时问题出现了：由 $x=y+tu_{0}(y)$ 可知，当 $x_{y}=1+tu_{0}'(y)$ 为零时变换奇异。

我们可以停在最早奇点出现的临界时刻，或者寻求所谓\emph{极大发展区域}，即上半平面上 $u$ 作为古典解存在的最大子集。在 Lagrange 坐标下，这是满足 $tu_{0}'(y)>-1$ 的 $(y,t)$ 集合。

为供后续参考，我们把\eqref{align:ch04-Burgers}解的行为与带阻尼的 Burgers 方程解的行为作一对比：
\begin{align}
\partial_{t}u(x,t)+\partial_{x}\left(\tfrac{1}{2}u^{2}(x,t)\right)+u(x,t)=0.
\label{align:ch04-Burgers阻尼}
\end{align}
把上面对\eqref{align:ch04-Burgers}的论证改造后用于\eqref{align:ch04-Burgers阻尼}，可得：沿特征线 $\mathrm{d}x/\mathrm{d}t=u(x,t)$，古典解 $u$ 及其导数 $v=\partial_{x}u$ 的演化现由方程 $\dot{u}+u=0$ 与 $\dot{v}+v^{2}+v=0$ 支配。后一方程例示了非线性的失稳作用与阻尼的平滑效应之间的竞争：当初值 $u_{0}$ 满足 $u_{0}'(x)\geqslant -1$（对所有 $x\in(-\infty,\infty)$）时，阻尼占优，$\partial_{x}u$ 保持有界，Cauchy 问题存在整体古典解。反之，若存在 $\bar{x}\in(-\infty,\infty)$ 使 $u_{0}'(\bar{x})<-1$，则波在有限时间内破碎，因为 $v=\partial_{x}u$ 必沿从点 $(\bar{x},0)$ 发出的特征线爆破。

为看到古典解崩溃的另一情形，考虑\eqref{align:ch04-Burgers}以支集在区间 $[0,1]$ 内的初值 $u_{0}(\cdot)$ 的 Cauchy 问题。设在某时间区间 $[0,T)$ 上存在古典解 $u(x,t)$。此时，由\cref{theorem:ch04-影响区域}，对任一 $t\in[0,T)$，$u(\cdot,t)$ 的支集在 $[0,1]$ 内。我们定义加权总质量
\begin{align}
M(t)=\int_{0}^{1}xu(x,t)\,\mathrm{d}x,
\label{align:ch04-加权质量}
\end{align}
并利用\eqref{align:ch04-Burgers}与 Schwarz 不等式导出微分不等式
\begin{align}
\dot{M}(t)=-\tfrac{1}{2}\int_{0}^{1}x\partial_{x}(u^{2}(x,t))\,\mathrm{d}x=\tfrac{1}{2}\int_{0}^{1}u^{2}(x,t)\,\mathrm{d}x\geqslant\tfrac{3}{2}M^{2}(t).
\label{align:ch04-M微分不等式}
\end{align}
于是，若 $M(0)>0$，则 $M(t)$ 必在时刻 $t^{*}=\tfrac{2}{3}M(0)^{-1}$ 或之前爆破。其解释是：\cref{theorem:ch04-影响区域}所施加的古典解影响范围增长速率的限制，把“质量”囚禁住，从而阻碍了色散。这导致质量正部与负部分离，并最终引起爆破。然而不难看出，在当前情形下波将在时刻 $\bar{t}=\tfrac{1}{6}M(0)^{-1}$ 之前开始破碎，即波破碎灾难将在质量囚禁灾难可能发生之前出现。

正如我们将在\cref{chapter:ch07-一维双曲平衡律组}中看到的，波破碎灾难对于单个空间变量中真正非线性的双曲守恒律组的解是普遍出现的，因为沿特征线传播的波被限制在一个平面内，从而不可避免地相互碰撞。相反，在多个空间维数中，波破碎可能因其与色散的竞争而被避免。尽管如此，正如我们将在\cref{chapter:ch06-标量守恒律L1理论}与本节第 4.8 小节中看到的，对于多个空间维数的标量守恒律以及三维 Euler 方程组，波仍然破碎。关于色散占优并阻止波破碎的方程组的讨论超出本书范围。一个来自弹性动力学的这类方程组的例子将在\cref{chapter:ch05-熵与经典解稳定性}中展示。


\end{document}
